% !TEX root = ../main.tex
% !TEX program = xelatex
\section{模型评价与推广}

\subsection{模型优点与创新}

\begin{enumerate}
    \item \textbf{结构统一，便于逐步扩展。}
    四问共用电量平衡与储能约束，逐步加入预测误差、日内调整和波动电价，既适应信息条件的变化，也便于解释各问联系和复核结果。

    \item \textbf{围绕实际费用处理预测偏差。}
    根据近期误差设置准备，并用日内负载修正判断，以适量普通购电减少昂贵补购。费用与风险共同衡量方案优劣，使预测服务于决策。

    \item \textbf{区分不同信息带来的收益。}
    同条件对照表明，日内反馈与改单节省50.63万元，后发光伏预报再节省6.95万元，明确了节费来源及新预报的引入价值。
\end{enumerate}

\subsection{模型局限与改进方向}

现有参数比较基于2025年数据，后续可通过跨年验证完善校准，补充供需与电价共同异常的情景，并纳入电池老化和网络损耗成本，使方案更贴近长期运行条件。

\subsection{模型推广性评估}

该框架可为校园、社区和工业园区微网提供参考。按当地供需、设备和交易规则校准后，可扩展至风光互补、多台储能及可调负载，协调本地发电与外网购电，兼顾用能经济性和可再生能源利用。
